\item ¿Por qué no es posible la siguiente situación? Un grupo de excursionistas se levanta a las 8:30 a.m. y utiliza una cocina solar que consiste en una superficie reflectora curva que concentra la luz solar sobre el objeto a calentar (figura P7.42). Durante el día, la máxima intensidad solar que llega a la superficie de la Tierra en la ubicación de la cocina es $I = 600\text{ W/m}^2$. La cocina encara al Sol y tiene un diámetro de cara $d = 0.600\text{ m}$. Suponga que una fracción $f$ de 40.0\% de la energía incidente se transfiere a 1.50 L de agua en un contenedor abierto, inicialmente a $20.0^\circ\text{C}$. El agua llega al punto de ebullición y los excursionistas disfrutan un café caliente en el desayuno antes de caminar 10 millas y retornar al mediodía para la comida.
